\section{Síntesis y ampliación}
\subsection{Tabla de síntesis principal}
\begin{tabularx}{\textwidth}{lXl}
\toprule
Dimensión & Idea central & Implicación práctica \\
\midrule
Señal de recompensa & Imágenes, lenguaje y success classifier sustituyen la reward manual & Al diseñar el reward pipeline, incorporar mecanismos de confianza/ensemble; \\
Operación del entorno & Forward-backward controllers y time-aware buffer hacen posible el aprendizaje reset-free & Cada reset policy necesita un registro de recovery y debe vincularse con el scheduler; \\
Gestión de objetivos & Goal-Query y curriculum policy logran una dificultad adaptativa & Mantener un goal repository y registrar novelty/skill dependency; \\
Seguridad y operación & safety penalty + regular monitoring evitan el drift & Construir un dashboard y conservar un canal de rollback; \\
\bottomrule
\end{tabularx}

\subsection{Lecturas adicionales}
\begin{itemize}
\item Sharma et al., ``Autonomous Reinforcement Learning: Formalism and Benchmarking,'' ICLR 2022
\item Eysenbach et al., ``Leave No Trace: Learning to Reset for Safe and Autonomous RL,'' ICLR 2018
\item Pong et al., ``Skew-Fit: State-Covering Self-Supervised RL,'' ICML 2020
\item Nair et al., ``Visual Reinforcement Learning with Imagined Goals,'' NeurIPS 2018
\item Ma et al., ``VIP: Towards Universal Visual Reward and Representation via Value-Implicit Pre-Training,'' ICLR 2023
\end{itemize}

\subsection{Resumen de la sección}
Mediante una exposición sistemática de cinco dimensiones (recompensa, reset, goal, seguridad y monitoreo), esta clase presentó la ruta completa del aprendizaje autónomo de robots y ofreció mecanismos ejecutables de retroalimentación y reversión para una Autonomy de nivel de despliegue.

\end{document}
